\documentclass[11pt,a4paper]{article}

\usepackage[a4paper,margin=1in]{geometry}
\usepackage{fontspec}
\IfFontExistsTF{Times New Roman}
  {\setmainfont{Times New Roman}}
  {\IfFontExistsTF{TeX Gyre Termes}{\setmainfont{TeX Gyre Termes}}{\setmainfont{Times}}}
\usepackage{newtxmath}
\usepackage{amsmath}
\usepackage{booktabs}
\usepackage{longtable}
\usepackage{array}
\usepackage{enumitem}
\usepackage{xcolor}
\usepackage{listings}
\usepackage{graphicx}
\usepackage{float}
\usepackage[round,authoryear]{natbib}
\usepackage[hidelinks]{hyperref}
\usepackage{xurl}
\usepackage{fancyhdr}
\usepackage{microtype}

\setlength{\parindent}{0pt}
\setlength{\parskip}{0.65em}
\setlength{\emergencystretch}{3em}
\setlist[itemize]{leftmargin=1.4em,itemsep=0.2em,topsep=0.2em}
\setlist[enumerate]{leftmargin=1.6em,itemsep=0.2em,topsep=0.2em}

\pagestyle{fancy}
\fancyhf{}
\lhead{\small Ticket 07 Study Report}
\rhead{\small Dividend-Paying American Call Validation}
\cfoot{\thepage}
\renewcommand{\headrulewidth}{0.3pt}

\lstset{
  basicstyle=\ttfamily\small,
  breaklines=true,
  frame=single,
  columns=fullflexible,
  keepspaces=true,
  rulecolor=\color{black!30},
  xleftmargin=0.5em,
  xrightmargin=0.5em
}

\newcommand{\code}[1]{\texttt{#1}}
\newcommand{\file}[1]{\path{#1}}
\newcommand{\func}[1]{\path{#1}}
\newcolumntype{L}[1]{>{\raggedright\arraybackslash}p{#1}}
\newcolumntype{P}[1]{>{\raggedright\arraybackslash\ttfamily\footnotesize}p{#1}}

\title{\textbf{Ticket 07 Study Report: Dividend-Paying American Call Validation}\\
\large SURF2026 American Option Risk Surfaces}
\author{Codex-assisted implementation report}
\date{June 16, 2026}

\begin{document}
\maketitle

\begin{abstract}
Ticket 07 validates the positive-dividend American call regime using the existing Ticket 04
American CN/PSOR solver without changing solver code. Unlike Ticket 06, where \(q=0\) and a
standard American call should match the European call, this ticket sets \(q=0.08\) and checks
whether the American solver can represent a regime where early exercise value may become
meaningful. The validation compares American call values against Ticket 01 European
Black-Scholes call values, checks the payoff obstacle \(U\geq\Phi\), records PSOR convergence
metadata, writes reproducible CSV tables, and produces three validation figures. This report is
written for beginner researchers with accounting, data, machine-learning, and Python experience
who are still learning American option pricing, dividends, obstacle problems, and CN/PSOR
validation.
\end{abstract}

\tableofcontents
\newpage

\section{Role of Ticket 07 in the Whole Solver Project}

The project is building a classical, validated American option benchmark before moving to
boundary diagnostics, Greeks, stress regimes, datasets, or neural surrogate models. Ticket 07 is
part of the validation ladder:
\begin{itemize}
  \item Ticket 01 created payoff functions and European Black-Scholes closed-form utilities.
  \item Ticket 02 created uniform finite-difference grids, spatial operator coefficients, and
        simple boundary value helpers.
  \item Ticket 03 validated European Crank-Nicolson time stepping against closed-form prices.
  \item Ticket 04 implemented the American CN/PSOR obstacle core.
  \item Ticket 05 validated representative American put behavior.
  \item Ticket 06 validated the no-dividend American call control theorem, \(q=0\).
\end{itemize}

Ticket 07 changes the financial regime from \(q=0\) to \(q>0\). It asks whether the American
exercise machinery can show meaningful American-minus-European value when dividends make early
exercise more plausible. It does not modify the solver. It uses the existing solver as a
black-box validation target and records whether the output is financially and numerically
credible.

\section{Theory and Methodology/Literature Connection}

The documented source basis for this report is internal:
\begin{itemize}
  \item \file{reports/01_literature/literature_map.md};
  \item \file{reports/02_math/formulation_note.md};
  \item \file{reports/03_solver/solver_validation_plan.md};
  \item \file{reports/03_solver/solver_implementation_tickets.md};
  \item the Ticket 01 through Ticket 06 reports and code.
\end{itemize}

The literature map and formulation note frame American options as obstacle problems or linear
complementarity problems when early exercise is allowed. Ticket 04 implemented that LCP logic
with PSOR. The solver validation plan then asks for staged validation: first European prices,
then American puts, then no-dividend American calls, then dividend-paying American calls. Ticket
07 follows that sequence.

This reference integration pass adds verified citations for the external theory behind the
project documents. Dividend-adjusted option pricing follows the Black-Scholes/Merton framework
\citep{black_scholes_1973,merton_1973}, and the American obstacle/LCP interpretation is connected
to standard complementarity methodology \citep{cottle_pang_stone_2009}.

\section{Why Dividends Can Make American Call Early Exercise Meaningful}

A call option gives the holder the right to buy the underlying stock for strike \(K\). The call
payoff is
\begin{equation}
  \Phi_{\mathrm{call}}(S) = \max(S-K,0).
\end{equation}

For a no-dividend stock, exercising a call early is usually unattractive. The holder gives up
remaining optionality and pays the strike early without receiving any dividend benefit. That is
why Ticket 06 checked the theorem
\begin{equation}
  C_{\mathrm{American}}(S,T;q=0) \approx C_{\mathrm{European}}(S,T;q=0).
\end{equation}

With a continuous dividend yield \(q>0\), the trade-off changes. Holding the option does not
receive the underlying dividend stream. Exercising the call gives ownership of the stock, and
stock ownership benefits from dividends. When the dividend benefit is large enough, especially
for sufficiently in-the-money calls, early exercise can become economically meaningful.

This does not mean every positive-dividend American call must exercise early under every
parameter set. It means the solver should be able to represent the possibility. Ticket 07
therefore checks for a positive American-minus-European value in a representative \(q=0.08\)
case, while still avoiding any free-boundary claim.

\section{Contrast with Ticket 06 No-Dividend Call Validation}

Ticket 06 was a control test: if \(q=0\), the American call should not be meaningfully above the
European call. Ticket 07 is the positive-dividend test: if \(q>0\), the American call should
still dominate payoff and the European call, and it may be meaningfully above the European call.

\begin{longtable}{L{0.23\textwidth}L{0.31\textwidth}L{0.34\textwidth}}
\toprule
Ticket & Financial setting & Main validation idea \\
\midrule
\endhead
Ticket 06 & No-dividend American call, \(q=0\) &
American value should approximately equal European value; no false early-exercise premium. \\
Ticket 07 & Dividend-paying American call, \(q=0.08\) &
American value should dominate payoff and European value; positive American-minus-European value
can appear. \\
\bottomrule
\end{longtable}

This contrast is important for beginner researchers. The same American solver can be tested in a
case where early exercise should not matter and then in a case where it can matter. Passing both
checks gives stronger evidence than looking at only one regime.

\section{Dividend-Paying American Call Financial Intuition}

For a European call with continuous dividends, the closed-form benchmark includes the dividend
discount factor \(e^{-qT}\). Intuitively, future stock value is reduced by the dividend yield in
the risk-neutral pricing formula:
\begin{equation}
  C_{\mathrm{European}}
  = S e^{-qT} N(d_1) - K e^{-rT} N(d_2).
\end{equation}

For an American call, the holder can also compare:
\begin{itemize}
  \item the continuation value: keep the option alive; and
  \item the immediate exercise value: take \(\max(S-K,0)\) now.
\end{itemize}

The American obstacle condition is
\begin{equation}
  U(S,\tau) \geq \Phi_{\mathrm{call}}(S).
\end{equation}
When \(U>\Phi\), continuation is better at that grid node. When \(U=\Phi\), immediate exercise is
competitive with continuation. Ticket 07 reports the obstacle gap \(U-\Phi\), but only as a
validation diagnostic. It does not infer an exercise boundary from that gap.

\section{Validation Criteria}

\subsection{Positive-Dividend Setup}

Every Ticket 07 case must use a call option and \(q>0\). The default cases use:
\[
  K=1,\quad T=1,\quad r=0.05,\quad q=0.08,\quad \sigma=0.20,\quad S_{\max}=4.
\]
The tests also verify that every validation and selected-spot row reports a positive \(q\).

\subsection{Obstacle Condition}

The American call value should never fall below payoff:
\begin{equation}
  U(S_i,\tau_n) - \Phi(S_i) \geq 0.
\end{equation}
The experiment reports the minimum obstacle gap and the maximum obstacle violation over the full
stored value grid. A zero maximum obstacle violation means the PSOR projection respected the
obstacle on the tested grids.

\subsection{American Versus European Call}

An American option gives at least as much exercise flexibility as the corresponding European
option. In the positive-dividend call case, the experiment checks:
\begin{equation}
  U_{\mathrm{American\ call}}(S,T;q>0)
  \geq
  C_{\mathrm{European\ call}}(S,T;q>0),
\end{equation}
up to numerical tolerance.

\subsection{Dividend-Call Early-Exercise Value Diagnostic}

The diagnostic difference is
\begin{equation}
  \Delta_{\mathrm{Am-Eur}}(S)
  = U_{\mathrm{American\ call}}(S,T) - C_{\mathrm{European\ call}}(S,T).
\end{equation}
Ticket 07 reports the minimum, maximum, maximum absolute difference, RMSE, and selected-spot
values. Positive \(\Delta_{\mathrm{Am-Eur}}\) is evidence that the solver is representing the
dividend-call early-exercise regime. It is not a boundary location.

\subsection{Target Region}

Summary error and difference metrics focus on:
\[
  0.4 \leq S/K \leq 1.8.
\]
This follows the existing solver validation plan. The selected-spot table additionally includes
\(S/K=2.0\) because the dividend-call effect is particularly visible for higher spot values.

\section{What Ticket 07 Is Not}

Ticket 07 does not redo the no-dividend call theorem validation from Ticket 06. It uses \(q>0\)
only.

Ticket 07 is also not boundary extraction. It does not compute a continuation premium module,
does not search for a free-boundary location, and does not label any figure as a boundary. The
American-minus-European curve is a pricing comparison, not a free-boundary diagnostic.

Finally, Ticket 07 does not implement Greeks, stress maps, datasets, or neural-network files.
Those require a more complete validated solver pipeline and belong to later work.

\section{Implementation Design Decisions}

\begin{longtable}{L{0.28\textwidth}L{0.60\textwidth}}
\toprule
Decision & Reason \\
\midrule
\endhead
Use the Ticket 04 solver unchanged &
Ticket 07 is a validation ticket. Changing the solver would mix implementation and validation
concerns. \\
Use \(q=0.08\) &
The dividend yield is positive and large enough to make American-minus-European value visible in
the representative case. \\
Use medium and fine grids &
The \(120\times120\) and \(180\times180\) cases give a light refinement audit without becoming a
full grid/domain sensitivity study. \\
Use nearest-grid selected spots &
Nearest-grid reporting is simple and reproducible for beginner readers. The report states that
this is not interpolated convergence analysis. \\
Generate three figures &
The value comparison, American-minus-European curve, and PSOR iteration curve each support a
specific validation question. \\
Avoid obstacle-gap figure &
Obstacle gaps are reported in tables only to reduce the chance of confusing this ticket with
boundary extraction. \\
\bottomrule
\end{longtable}

\section{Code-Level Function Explanations}

\begin{longtable}{L{0.38\textwidth}L{0.50\textwidth}}
\toprule
Function & Role \\
\midrule
\endhead
\func{run_validation_cases} &
Runs the medium and fine positive-dividend American call cases, then builds validation rows,
selected-spot rows, and figure metadata. \\
\func{summarize_dividend_call_case} &
Computes American-versus-European metrics, obstacle diagnostics, positive-difference node count,
and PSOR convergence summaries for one solver result. \\
\func{selected_spot_rows} &
Reports payoff, European call, American call, American-minus-European, and American-minus-payoff
at selected moneyness points. \\
\func{nearest_spot_index} &
Chooses the nearest grid node to a requested spot so selected-spot reporting is reproducible. \\
\func{write_csv} &
Writes stable CSV artifacts with fixed column ordering. \\
\func{create_value_comparison_figure} &
Creates the American call, European call, and payoff comparison figure if matplotlib is
available. \\
\func{create_american_minus_european_figure} &
Creates the American-minus-European curve for the dividend-call value diagnostic. \\
\func{create_psor_iterations_figure} &
Creates the PSOR iteration-count curve by time step to summarize convergence behavior. \\
\code{main} &
Runs the full Ticket 07 experiment and writes CSV and figure artifacts. \\
\bottomrule
\end{longtable}

\section{Validation Result Tables and Interpretation}

\subsection{Validation Cases and Parameters}

\begin{longtable}{@{}L{0.14\textwidth}L{0.07\textwidth}L{0.07\textwidth}L{0.07\textwidth}L{0.07\textwidth}L{0.09\textwidth}L{0.07\textwidth}L{0.07\textwidth}L{0.08\textwidth}@{}}
\toprule
Case & \(K\) & \(T\) & \(r\) & \(q\) & \(\sigma\) & \(M\) & \(N\) & \(S_{\max}\) \\
\midrule
\endhead
Medium & 1.0 & 1.0 & 0.05 & 0.08 & 0.20 & 120 & 120 & 4.0 \\
Fine & 1.0 & 1.0 & 0.05 & 0.08 & 0.20 & 180 & 180 & 4.0 \\
\bottomrule
\end{longtable}

Both validation cases use \(q=0.08>0\). This distinguishes Ticket 07 from the no-dividend
Ticket 06 control case.

\subsection{American-Versus-European Summary}

{\footnotesize
\setlength{\tabcolsep}{3pt}
\begin{longtable}{@{}L{0.09\textwidth}L{0.13\textwidth}L{0.11\textwidth}L{0.12\textwidth}L{0.12\textwidth}L{0.10\textwidth}L{0.10\textwidth}@{}}
\toprule
Case & Max abs diff & RMSE & Min Am. - Eur. & Max Am. - Eur. & Max diff spot & Positive nodes \\
\midrule
\endhead
Medium & \(8.9422\times10^{-2}\) & \(4.2654\times10^{-2}\) & \(4.1188\times10^{-8}\) & \(8.9422\times10^{-2}\) & 1.8 & 43 \\
Fine & \(8.9422\times10^{-2}\) & \(4.2454\times10^{-2}\) & \(1.6879\times10^{-8}\) & \(8.9422\times10^{-2}\) & 1.8 & 64 \\
\bottomrule
\end{longtable}
}

The American-minus-European difference is nonnegative in the target region within the reported
tolerance and is materially positive in the representative dividend case. The maximum target
region difference is about \(0.0894\), and the selected \(S/K=2.0\) row shows a difference of
about \(0.1050\). This is the key financial evidence in Ticket 07: the solver does not collapse
the positive-dividend call case into the no-dividend theorem behavior.

\subsection{Obstacle Summary}

\begin{longtable}{@{}L{0.16\textwidth}L{0.18\textwidth}L{0.22\textwidth}@{}}
\toprule
Case & Min obstacle gap & Max obstacle violation \\
\midrule
\endhead
Medium & 0 & 0 \\
Fine & 0 & 0 \\
\bottomrule
\end{longtable}

The obstacle diagnostics show no violation in the stored value grids. This confirms that the
American call values were projected above payoff during PSOR time stepping. Obstacle validity is
necessary, but it is not sufficient by itself. Ticket 07 also needs the European comparison and
PSOR convergence evidence.

\subsection{PSOR Convergence Metadata}

{\footnotesize
\setlength{\tabcolsep}{3pt}
\begin{longtable}{@{}L{0.10\textwidth}L{0.17\textwidth}L{0.11\textwidth}L{0.13\textwidth}L{0.13\textwidth}L{0.16\textwidth}@{}}
\toprule
Case & All steps converged & Step count & Max iterations & Mean iterations & Max final update \\
\midrule
\endhead
Medium & True & 120 & 10 & 7.35 & \(9.4309\times10^{-9}\) \\
Fine & True & 180 & 11 & 7.1111 & \(9.9674\times10^{-9}\) \\
\bottomrule
\end{longtable}
}

Every PSOR step converged in both cases. The maximum final update remained below the intended
\(10^{-8}\) tolerance scale. This makes the American-versus-European comparison more credible:
the reported dividend-call premium is not being produced by visible PSOR non-convergence.

\subsection{Selected Spot Price Comparison}

{\footnotesize
\begin{longtable}{@{}L{0.10\textwidth}L{0.09\textwidth}L{0.11\textwidth}L{0.13\textwidth}L{0.13\textwidth}L{0.13\textwidth}L{0.13\textwidth}@{}}
\toprule
Case & \(S/K\) & Payoff & European call & American call & Am. - Eur. & Am. - payoff \\
\midrule
\endhead
Medium & 0.5 & 0 & 0.00000484 & 0.00000633 & 0.00000149 & 0.00000633 \\
Medium & 0.8 & 0 & 0.00817709 & 0.00836716 & 0.00019007 & 0.00836716 \\
Medium & 1.0 & 0 & 0.06142998 & 0.06513118 & 0.00370119 & 0.06513118 \\
Medium & 1.2 & 0.20000000 & 0.18282650 & 0.20369502 & 0.02086852 & 0.00369502 \\
Medium & 1.5 & 0.50000000 & 0.43612545 & 0.50000000 & 0.06387455 & 0 \\
Medium & 2.0 & 1.00000000 & 0.89503490 & 1.00000000 & 0.10496510 & 0 \\
Fine & 0.488889 & 0 & 0.00000301 & 0.00000352 & 0.00000051 & 0.00000352 \\
Fine & 0.8 & 0 & 0.00817709 & 0.00842129 & 0.00024420 & 0.00842129 \\
Fine & 1.0 & 0 & 0.06142998 & 0.06529167 & 0.00386168 & 0.06529167 \\
Fine & 1.2 & 0.20000000 & 0.18282650 & 0.20376496 & 0.02093846 & 0.00376496 \\
Fine & 1.488889 & 0.48888889 & 0.42612581 & 0.48888889 & 0.06276307 & 0 \\
Fine & 2.0 & 1.00000000 & 0.89503490 & 1.00000000 & 0.10496510 & 0 \\
\bottomrule
\end{longtable}
}

The selected-spot rows make the dividend-call effect readable. Near \(S/K=1.0\), the American
call is modestly above the European call. At \(S/K=1.5\) and \(S/K=2.0\), the American value is at
payoff in the final-time slice and is meaningfully above the European value. The report does not
use these rows to infer an exercise boundary; they are selected price checks only.

\subsection{Medium-to-Fine Comparison}

The largest selected-spot medium-to-fine value difference reported by the experiment is
\(0.0111111\). This is a nearest-grid audit, not an interpolated convergence proof. It is driven
mainly by selected spots where one grid lands exactly on a payoff node and the other grid lands on
the nearest neighboring node. A full grid and domain sensitivity study belongs to a later ticket.

\section{Figure Interpretation}

Figure~\ref{fig:value-comparison} plots the fine-grid American call value, European call value,
and payoff over moneyness. The American curve is above the European curve in the dividend-paying
case, especially for higher spot values. The payoff curve provides the lower obstacle that the
American value must respect.

\begin{figure}[H]
\centering
\includegraphics[width=0.82\textwidth]{results/01_solver_validation/figures/ticket_07_dividend_american_call_value_comparison.png}
\caption{Ticket 07 dividend-paying American call value, European call value, and payoff comparison.}
\label{fig:value-comparison}
\end{figure}

Figure~\ref{fig:am-eur} plots American call value minus European call value. This directly shows
the positive-dividend value diagnostic. The zero line is included for reference. No boundary is
extracted or labeled.

\begin{figure}[H]
\centering
\includegraphics[width=0.82\textwidth]{results/01_solver_validation/figures/ticket_07_dividend_american_call_american_minus_european.png}
\caption{Ticket 07 American-minus-European value diagnostic for the dividend-paying call.}
\label{fig:am-eur}
\end{figure}

Figure~\ref{fig:psor-iterations} shows the number of PSOR iterations used at each time step in
the fine case. The curve is a convergence-behavior diagnostic: it helps show that the validation
prices come from converged LCP solves, not hidden failed iterations.

\begin{figure}[H]
\centering
\includegraphics[width=0.82\textwidth]{results/01_solver_validation/figures/ticket_07_dividend_american_call_psor_iterations.png}
\caption{Ticket 07 PSOR iteration counts by time step for the fine positive-dividend call case.}
\label{fig:psor-iterations}
\end{figure}

\section{Testing Strategy}

\begin{longtable}{L{0.27\textwidth}L{0.31\textwidth}L{0.30\textwidth}}
\toprule
Test category & What it proves & What it does not prove \\
\midrule
\endhead
Script/artifact tests &
The experiment runs and writes the expected Ticket 07 CSV and PNG artifacts. &
Full validation across all dividend regimes. \\
CSV schema tests &
The validation and selected-spot tables have stable columns. &
That every later diagnostic has already been designed. \\
\(q>0\) tests &
Every validation row is genuinely positive-dividend. &
Sensitivity to different dividend yields. \\
Obstacle tests &
American call values dominate payoff within tolerance. &
That a free boundary has been found. \\
European comparison tests &
American values dominate European values within tolerance. &
A formal convergence proof. \\
Positive-difference tests &
The representative dividend case shows visible American-minus-European value. &
That all positive-dividend calls must exercise early. \\
PSOR metadata tests &
Default time steps converge and expose iteration/update summaries. &
Convergence under every possible grid or relaxation parameter. \\
Nearest-grid test &
Selected-spot reporting is predictable. &
Interpolation accuracy. \\
Figure tests &
Meaningful PNGs are created when matplotlib is available. &
That plots alone validate the method. \\
Scope and docstring tests &
Ticket 07 files identify their ticket and do not introduce forbidden APIs/files. &
Future tickets that intentionally add diagnostics. \\
\bottomrule
\end{longtable}

\section{Test Results and Interpretation}

The final test command was:
\begin{lstlisting}[language=bash]
PYTHONPYCACHEPREFIX=/private/tmp/codex_pycache PYTHONPATH=src MPLCONFIGDIR=/private/tmp/matplotlib-cache FC_CACHEDIR=/private/tmp/fontconfig-cache python3 -m unittest discover -s tests
\end{lstlisting}

Observed result during Ticket 07 implementation:
\begin{lstlisting}
Ran 76 tests

OK
\end{lstlisting}

The experiment command was:
\begin{lstlisting}[language=bash]
PYTHONPYCACHEPREFIX=/private/tmp/codex_pycache PYTHONPATH=src MPLCONFIGDIR=/private/tmp/matplotlib-cache FC_CACHEDIR=/private/tmp/fontconfig-cache python3 experiments/04_dividend_american_call_validation.py
\end{lstlisting}

The experiment wrote two validation rows, twelve selected-spot rows, and three PNG figures.

\section{Implementation Pitfalls and Debugging Notes}

\begin{longtable}{@{}L{0.22\textwidth}L{0.24\textwidth}L{0.22\textwidth}L{0.18\textwidth}@{}}
\toprule
Recorded item or pitfall & Symptom & Fix or guard & Beginner lesson \\
\midrule
\endhead
Expected TDD red run &
The new tests failed with \code{FileNotFoundError} because the Ticket 07 experiment script did
not exist yet. &
Create \file{experiments/04_dividend_american_call_validation.py}. &
A red test confirms the test is aimed at the new artifact. \\
Accidentally using \(q=0\) &
The experiment would repeat Ticket 06 rather than testing the dividend regime. &
Validate every case and output row has \(q>0\). &
The theorem regime is controlled by assumptions, not by filenames. \\
Overclaiming dividend behavior &
A positive dividend yield does not guarantee early exercise under every parameter set. &
Phrase results as evidence for this representative \(q=0.08\) case. &
Validation evidence is parameter-specific unless a broader study is run. \\
Confusing premium diagnostics with boundary extraction &
American-minus-European and obstacle gap curves can tempt readers to infer an exercise boundary. &
Report pricing differences and state that no boundary is extracted or labeled. &
Diagnostics and boundary algorithms are different tasks. \\
Checking only \(U\geq\Phi\) &
The solver could respect payoff but still fail the European comparison. &
Include American-versus-European rows, metrics, and figures. &
American validation needs multiple consistency checks. \\
Ignoring PSOR metadata &
Positive premiums could be suspect if time steps did not converge. &
Report convergence flags, iteration counts, and final update values. &
Numerical evidence should include solver-status evidence. \\
Creating too many figures &
Extra figures can distract from the validation question. &
Use three figures, each with a distinct validation purpose. &
Every plot should answer a specific question. \\
\bottomrule
\end{longtable}

No unexpected solver bugs were recorded during this ticket. The only failed test iteration was
the planned red TDD run before the Ticket 07 experiment script existed.

\section{Limitations}

\begin{longtable}{L{0.34\textwidth}L{0.54\textwidth}}
\toprule
Not included in Ticket 07 & Reason \\
\midrule
\endhead
No-dividend call validation & Already covered in Ticket 06. \\
American put validation & Already covered in Ticket 05. \\
Free-boundary extraction & Ticket 07 does not infer or label exercise boundaries. \\
Continuation-premium module & Boundary-oriented diagnostics are later work. \\
Greeks & Delta and Gamma diagnostics are later work. \\
Stress maps, datasets, neural networks & These require a more complete validated solver pipeline. \\
Full grid/domain sensitivity & The current medium-to-fine comparison is preliminary only. \\
Broad dividend sensitivity & The ticket uses one representative positive-dividend setup, \(q=0.08\). \\
\bottomrule
\end{longtable}

\section{Lessons Learned / Study Notes}

The first lesson is that dividend yield changes American call economics. With \(q=0\), early
exercise should not be valuable for a standard call. With \(q>0\), early exercise can become
meaningful because stock ownership receives dividend benefits that option ownership does not.

The second lesson is that American-minus-European value is a pricing diagnostic, not
automatically a boundary. It can show that the American exercise right has value, but boundary
extraction requires separate logic and careful definitions.

The third lesson is that \(U\geq\Phi\), American \(\geq\) European, and PSOR convergence answer
different questions. A credible validation ticket should check all three.

The fourth lesson is that selected-spot tables are for human inspection. They help students see
what is happening at familiar moneyness levels, but they do not replace full convergence or
sensitivity analysis.

\section{Link to the Next Ticket}

Ticket 07 completes the basic American call validation pair:
\begin{itemize}
  \item Ticket 06 showed the no-dividend control behavior: American call approximately equals
        European call when \(q=0\).
  \item Ticket 07 showed positive-dividend behavior: American call can be meaningfully above the
        European call while respecting payoff and PSOR convergence checks.
\end{itemize}

The next ticket can build on this validated solver ladder to add a new diagnostic layer. Later
work still needs boundary extraction, complementarity diagnostics, Greeks, grid/domain
sensitivity, stress maps, datasets, and neural surrogate work. Those are deliberately outside
Ticket 07.

\section{Reproducibility Record}

\begin{longtable}{L{0.27\textwidth}L{0.61\textwidth}}
\toprule
Item & Record \\
\midrule
\endhead
Experiment script &
\file{experiments/04_dividend_american_call_validation.py}. \\
Tests &
\file{tests/test_dividend_american_call_validation.py}. \\
Validation CSV &
\file{results/01_solver_validation/tables/ticket_07_dividend_american_call_validation.csv}. \\
Selected-spots CSV &
\file{results/01_solver_validation/tables/ticket_07_dividend_american_call_selected_spots.csv}. \\
Figures &
\file{results/01_solver_validation/figures/ticket_07_dividend_american_call_value_comparison.png};
\file{results/01_solver_validation/figures/ticket_07_dividend_american_call_american_minus_european.png};
\file{results/01_solver_validation/figures/ticket_07_dividend_american_call_psor_iterations.png}. \\
Report source &
\file{reports/03_solver/tickets/ticket_07_dividend_american_call_validation.tex}. \\
Report PDF &
\file{reports/03_solver/tickets/ticket_07_dividend_american_call_validation.pdf}. \\
\bottomrule
\end{longtable}

Commands used:
\begin{lstlisting}[language=bash]
PYTHONPYCACHEPREFIX=/private/tmp/codex_pycache PYTHONPATH=src MPLCONFIGDIR=/private/tmp/matplotlib-cache FC_CACHEDIR=/private/tmp/fontconfig-cache python3 experiments/04_dividend_american_call_validation.py
PYTHONPYCACHEPREFIX=/private/tmp/codex_pycache PYTHONPATH=src MPLCONFIGDIR=/private/tmp/matplotlib-cache FC_CACHEDIR=/private/tmp/fontconfig-cache python3 -m unittest discover -s tests
PYTHONPYCACHEPREFIX=/private/tmp/codex_pycache PYTHONPATH=src python3 -m compileall -q src tests experiments
\end{lstlisting}

\bibliographystyle{plainnat}
\bibliography{reports/03_solver/references}

\end{document}
